\chapter{Tecnologías de Back-End}
\pagestyle{fancy}

    Para el desarrollo de Back-End es ideal escoger la arquitectura y, sobre todo, el lenguaje y el Framework con el que se va a realizar el código fuente. Los lenguajes más populares para el desarrollo de Back-End descritos a continuación, fueron seleccionados de acuerdo a múltiples encuestas y estadísticas como 
    \href{https://survey.stackoverflow.co/2022/#technology-most-popular-technologies}{\textit{Stack Overflow Most Popular Technologies 2022}}, \href{https://www.jetbrains.com/lp/devecosystem-2022/}{\textit{2022 State of Developer Ecosystem Report by JetBrains}} y \href{https://www.tiobe.com/tiobe-index/}{\textit{TIOBE Index}}.

    \begin{figure}[!htbp]
        \centering
        \includegraphics[width=12cm]{graphics/teg-c6/TIOBE_INDEX_LANGUAGES.png}
        \caption{Lenguajes más populares en los últimos 10 años. TIOBE Index}
        \label{fig:tiobe_index_languages}
    \end{figure}

    Es posible determinar, gracias a la Fig. \ref{fig:tiobe_index_languages} y los otros dos indicadores previamente descritos, los 10 lenguajes más populares en el mercado en los últimos 10 años, sobre los que destacan los ampliamente conocidos lenguajes para el desarrollo Web: \textbf{JavaScript}, \textbf{Python} y \textbf{Java}. Es importante considerar que se toman en consideración no únicamente para el ámbito Web, sino para el desarrollo de aplicaciones en general. El índice de TIOBE provee dichos datos gracias a su metodología que involucra el análisis de las consultas realizadas en los motores de búsqueda relacionadas con lenguajes de programación, que da indicadores de la frecuencia con la que los desarrolladores están buscando por recursos para un lenguaje particular. Además de eso, toma otros factores más complejos como el número de ingenieros capacitados, cursos y programas de terceros para un lenguaje. Finalmente, también toma en cuenta la cantidad de proyectos de código abierto y la oferta laboral para cada lenguaje en su lista.

    \section[Micro-servicios]{Tecnologías para Micro-Servicios}
    
    El lenguaje de programación que se utilice para la implementación de micro-servicios debe tener las siguientes características claves

    \begin{enumerate}
        \item \textbf{Soporte para computación distribuida}: Los micro-servicios están típicamente distribuidos a través de distintos servidores e instancias, por lo que el lenguaje debería tener soporte para mensajería y computación distribuida, es decir, permitir la resolución de un problema a través del uso y comunicación de múltiples sistemas interconectados.
        \item \textbf{Ligero y rápido}: Los micro-servicios están pensados para ser componentes independientes y ligeros, por lo que el lenguaje de programación debería tener también estas características para minimizar el \textit{Overhead}, que se refiere al costo, en términos de recursos, de levantar un micro-servicio comparado con el costo de las operaciones que realiza. Es decir, no es rentable un micro-servicio si el costo de levantarlo es más alto o similar al costo computacional de las tareas que realiza.
        \item \textbf{Escalabilidad}: Una arquitectura de micro-servicios está diseñada para ser altamente escalable, por lo que el lenguaje que se escoja debe ser capaz de manejar grandes volúmenes de solicitudes y ser fácilmente escalable horizontalmente a través de múltiples servidores.
        \item \textbf{Resiliencia y tolerancia al fallo}: Los micro-servicios necesitan ser resilientes a los errores y fallos, por lo que el lenguaje de programación debería soportar patrones de diseño tolerantes al fallo, como \textit{Circuit-Breakers}.
        \item \textbf{Containerización}: Comúnmente los micro-servicios son desplegados usando tecnologías de \textit{containerización} como Docker, por lo que el lenguaje de programación debe poder soportarlo y ser compatible con estas plataformas.
    \end{enumerate}

Además de los lenguajes,  para una implementación efectiva de un ambiente de micro-servicios, suele ser común el uso de herramientas adicionales que ayudan a la creación de un ecosistema de servicios efectivo, como es el uso de Docker, Kubernetes, AWS, Jenkins, Zipkin, entre otros.

    \section{JavaScript}
        Es un lenguaje de programación interpretado, orientado a objetos y de alto nivel. Comúnmente se utiliza para construir aplicaciones Web y Scripts del lado del Cliente. Fue creado en 1995 por Brendan Eich y desde entonces se ha convertido en uno de los lenguajes más utilizados en el mundo.

        JavaScript se usa principalmente para crear Scripts del lado del Cliente en los navegadores Web, sin embargo, es posible utilizarlo también del lado del Servidor gracias a plataformas como Node.js, que permite el desarrollo de aplicaciones en este lenguaje fuera de un navegador Web. Este lenguaje provee múltiples características, como programación funcional, programación orientada a objetos e incluso programación orientada a eventos. También soporta programación asíncrona.

        JavaScript posee una larga y activa comunidad, con muchas librerías y Frameworks disponibles para desarrollo Web, pero también para todo tipo de Software. \cite{flanagan2011javascript}

        %% REVISAR SUBSECTIONS %%
        
        \subsection{Express.js}
            Es un Framework para Node.js de Back-End muy popular y ampliamente utilizado en el desarrollo Web. Provee una API simple y flexible para el desarrollo de aplicaciones Web, como también de APIs. Soporta una variedad de características como lo son enrutamientos, Middleware, y motores de plantillas. 
    
            Express.js tiene una comunidad grande y activa, y es ampliamente utilizado por aplicaciones en producción. \cite{expressjs}
        \subsection{NestJS}
            Aunque relativamente nuevo, es un Framework que está creciendo muy rápido y se utiliza para el desarrollo del Back-End de aplicaciones Web para JavaScript. Provee una arquitectura modular y escalable que soporta características de grandes Frameworks como inyección de dependencias, micro-servicios, GraphQL, entre otros.
            
            NestJS ha ganado popularidad en los años recientes, particularmente en el sector empresarial para aplicaciones de gran escala. \cite{nestjs}
    \section{Python}
        Python es una lenguaje de alto nivel e interpretado, lanzado por primera vez en 1991. Fue creado por Guido van Rossum y desde entonces se ha convertido en uno de los lenguajes de programación más populares en el mundo. Python es conocido por su simplicidad y legibilidad, lo que lo convierte en un lenguaje ideal tanto para principiantes como para desarrolladores experimentados.

        Es un lenguaje tipado dinámicamente, lo que significa que sus variables no necesitan ser declaradas con un tipo de dato específico. Posee una gran cantidad de librerías, que permiten desde el desarrollo Web hasta su uso para el análisis de datos en la investigación científica.

        Python soporta múltiples paradigmas de programación, como el procedimental, funcional y  orientado a objetos. Posee una comunidad grande y activa, y es el lenguaje número uno para Machine Learning e Inteligencia Artificial. \cite{downey2015think}
        \subsection{Django}
            Es un Framework Web de alto nivel lanzado en el año 2005. Concebido para alentar desarrollos rápidos y limpios y diseños pragmáticos, se ha convertido en una opción popular para el desarrollo de aplicaciones Web complejas y orientadas a los datos. Django provee un conjunto de herramientas como lo son un ORM, para fácil interacción con Bases de Datos, sistema de autenticación, interfaz de administrador, entre otros, que permiten a los desarrolladores levantar una aplicación funcional rápidamente.

            Una de las fortalezas de Django es su énfasis en la convención sobre la configuración. Django provee configuraciones por defecto basadas en mejores prácticas para que los desarrolladores se enfoquen en construir la aplicación sobre configurar el Framework. \cite{django}
        \subsection{Flask}
            Es un Framework Web ligero lanzado en el año 2010. Es conocido por su simplicidad y flexibilidad, lo que lo convierte en una opción popular para construir aplicaciones Web pequeñas y medianas.

            Una de las fortalezas principales de Flask es su minimalismo. Provee sólo lo esencial para construir una aplicación Web, como enrutamiento, motor de plantillas y manejo de peticiones, dejando el resto en manos de los desarrolladores para que puedan escoger las herramientas y librerías de su preferencia. 

            Flask posee una comunidad grande y activa, que ha contribuido en la creación de librerías que se integran fácilmente con el Framework \cite{flask}
    \section{Java}
        Java es un lenguaje de Programación muy popular, de propósito general, lanzado por primera vez en el año 1995 por la empresa Sun Microsystems. Diseñado para ser portátil, seguro y fácil de aprender, se convirtió rápidamente en una opción para desarrollar desde aplicaciones de escritorio, móviles hasta sistemas empresariales de gran escala. 

        Una de las características principales de Java es su habilidad de correr en múltiples plataformas, esto gracias a la Máquina Virtual de Java (JVM, por sus siglas en inglés). La JVM permite a los desarrolladores escribir el código una vez y que sea pueda ejecutar en cualquier plataforma que soporta la JVM, lo que representa un ahorro considerable de tiempo y esfuerzo al desarrollar y desplegar aplicaciones.

        Una gran fortaleza de Java es un simpleza y su legibilidad. El lenguaje es fácil de entender y escribir, lo que lo hace una gran opción para todo tipo de desarrolladores. La sintaxis de Java es limpia y consistente, lo que ayuda a reducir errores y mejorar la mantenibilidad del código. El sistema fuertemente tipado y la administración de memoria de Java lo hacen un lenguaje seguro y confiable para construir aplicaciones robustas.

        En general, la combinación que ofrece Java de simplicidad, legibilidad y versatilidad lo han convertido en uno de los lenguajes de programación más utilizados a nivel mundial. \cite{sierra2018headfirst}
        
        \subsection{Spring Boot}
            Es un Framework de código abierto poderoso y ampliamente utilizado para construir aplicaciones Web. Provee un conjunto de herramientas y librerías para crear aplicaciones de alto nivel, con una configuración inicial mínima. Spring Boot ofrece características, incluyendo la auto-configuración, servidores embebidos y fácil integración con otros proyectos Spring o librerías populares de terceros.

            Una de las características claves de Spring Boot es su facilidad de uso y desarrollo eficaz. Provee un conjunto de dependencias y configuraciones pre-definidas que permiten a los desarrolladores desplegar sus aplicaciones rápidamente. Además, Spring Boot ofrece una amplia documentación y soporte de la comunidad, convirtiéndolo en una opción muy popular para desarrolladores de todos los niveles de experiencia. \cite{wallis2019springboot}
        \subsection{Apache Struts}
            es un Framework de código abierto para construir aplicaciones Web en Java. Está basado en el patrón de Arquitectura de Modelo-Vista-Controlador (MVC) provee un conjunto de librerías para construir aplicaciones Web escalables y mantenibles. Struts ayuda a los desarrolladores a separar la capa de presentación (o vistas) de la lógica de negocio (el modelo) y la entrada del usuario y sus interacciones (el controlador).

            Una de las características claves de Struts es el soporte para crear aplicaciones Web basadas en formularios. Proporciona un conjunto de librerías y herramientas de validación que ayudan a simplificar el proceso de construir y validar formularios HTML. Además, Struts da soporte a un gran rango de tipos de entradas, incluyendo subida de archivos y seleccionadores de fechas e integra bien con otras librerías y Frameworks populares en el ecosistema Java. \cite{brown2019struts}
    